% Generated draft: requires mathematical and source review.
\documentclass[11pt]{article}
\usepackage[margin=25mm]{geometry}
\usepackage{fontspec}
\setmainfont{Latin Modern Roman}
\usepackage{amsmath,amssymb,mathtools,unicode-math}
\setmathfont{Latin Modern Math}
\usepackage{enumitem,needspace,xurl,hyperref}
\setlength{\parindent}{0pt}
\setlength{\parskip}{5pt}
\setlength{\emergencystretch}{3em}
\newcommand{\N}{\mathbb{N}}
\newcommand{\Z}{\mathbb{Z}}
\newcommand{\Q}{\mathbb{Q}}
\newcommand{\R}{\mathbb{R}}
\newcommand{\C}{\mathbb{C}}
\newcommand{\F}{\mathbb{F}}
\newcommand{\A}{\mathbb{A}}
\newcommand{\PP}{\mathbb{P}}
\newcommand{\E}{\mathbb{E}}
\newcommand{\eps}{\varepsilon}
\newcommand{\dd}{\,\mathrm{d}}
\begin{document}
\section*{UnsolvedMath ID: 30000707}
\textbf{Source problem code:} catalogue entry
\subsection{Definitions and mathematical statement}
In the usual entire-plane interpretation, let \(f,g\) be meromorphic on \(\C\), sharing four values \(a_1,a_2,a_3,\infty\), with the first three finite and distinct under the customary four-value convention. The source does not state whether sharing is with multiplicities (CM) or only equality of inverse-image sets (IM); these are different hypotheses. At \(\infty\), IM means equality of the pole sets. No multiplicity convention is chosen silently.

Retain exactly the restored auxiliary functions
\[
 \psi_1=\frac{f'(f-g)}{\prod_{\nu=1}^3(f-a_\nu)},\qquad
 \psi_2=\frac{g'(f-g)}{\prod_{\nu=1}^3(g-a_\nu)}.
\]
They are interpreted as meromorphic quotients, with removable singularities resolved where applicable, and the equations below are meromorphic identities. Valid solutions necessarily avoid identically undefined denominators; the nonzero right-hand sides exclude constant functions or \(f\equiv g\). The sign of \(f-g\) in the second quotient is the one printed in the source and is not changed to \(g-f\).

The two separate classification problems are
\[
 (\psi_1,\psi_2)=(e^z,e^{-z})
 \quad\text{and}\quad
 (\psi_1,\psi_2)=(e^{z^2},e^{-z^2}).
\]
For each system determine all qualifying pairs, up to the stated simultaneous transformations
\[
 (f(z),g(z))\longmapsto
 \bigl(T(f(az+b)),T(g(az+b))\bigr),
 \qquad a\in\C^\times,\quad b\in\C,
\]
where \(T(w)=(\alpha w+\beta)/(\gamma w+\delta)\) with \(\alpha\delta-\beta\gamma\ne0\). The same \(T,a,b\) must be used for both functions. Classification includes completeness and the possibility of an empty solution set; no proposed example or nonexistence claim is asserted as established.

The transformation convention needs care. Shared values change under \(T\), and the value normalized to \(\infty\) may cease to be infinite. Affine precomposition also changes derivatives by a factor and the written exponential arguments. Thus arbitrary displayed transformations need not preserve either normalized system with fixed \(a_\nu\) and exact right-hand sides. One possible equivalence is intersection of ambient transformation orbits with the normalized solution sets; an alternative permits re-normalizing shared values and auxiliary constants. The source does not specify which, so it is not claimed that the entire transformation family acts on a fixed normalized system.

The domain, ordering and fixed-versus-variable status of the shared values are not fully specified in the record. The literature triage checked 2026-08-21 reports a proposed first-system example and suspected second-system nonexistence, but says these are not a complete published resolution. Neither the example nor a proof is supplied. The open label and direct-source restoration note remain dated metadata.
\subsection{Short English statement}
Classify both normalized exponential auxiliary systems for four-value-sharing meromorphic pairs, using the same Möbius and affine changes on both functions. The sharing multiplicities and normalization of this equivalence are unspecified. The source's suggested example and suspected nonexistence remain unproved reports, not conclusions of the draft.
\subsection{Sources}
\begin{itemize}
\item Local UnsolvedMath registry, permanent ID 30000707, source code catalogue entry, ``Exponential Auxiliary Systems for Four-Value-Sharing Meromorphic Functions,'' including both restored quotients, both systems and the stated equivalence.
\item \emph{Normal Families and Complex Dynamics (2007)}, Oberwolfach Reports, pp.\ 487--548, DOI 10.4171/owr/2007/09: \url{https://doi.org/10.4171/owr/2007/09}.
\item Supplied reference ``four-value uniqueness background'': \url{https://doi.org/10.1112/jlms/s2-20.3.457}.
\item Literature triage checked 2026-08-21; record updated 2026-09-16. No external source has been visited or independently checked here.
\end{itemize}
\end{document}
